\section{Introduction}
\label{sec:intro}

Bid-rigging cartels in public procurement face a practical constraint:
they must simulate competition to avoid detection. A common solution is
\emph{cover bidding}---the deployment of firms whose role is not to win
auctions but to submit losing bids that legitimize the procurement
process, satisfy minimum-bidder requirements, and create the appearance
of competitive participation
\citep{marshall2012economics, asker2010leniency}. In procurement data,
these firms appear as bidders that repeatedly participate but never
win---what we call \emph{frequent losers} (FL). This paper documents
that FL presence is a robust marker of price anomalies consistent with
cartel activity, and provides suggestive evidence that cover bidding
is the underlying mechanism.

Using 4.5 million tender-items and 40 million bids from S\~ao Paulo's
Bolsa Eletr\^onica de Compras (BEC) platform over 2009--2019, we
identify 2,735 FL firms and document that tenders in which they
participate exhibit significantly higher prices. OLS with item, year,
and purchasing-unit fixed effects yields a 6.4\% price premium;
temporal cross-fitting---which defines FL status on one half of the
sample and estimates the price equation on the other, breaking any
mechanical link between classification and outcome---yields 3.6\%.
Matching estimators (CEM: 7.7\%; IPW: 5.5\%) bracket these figures.
Conditional on item fixed effects, the FL--non-FL price gap narrows
from 2.43 to 0.050 log points (overlap coefficient 0.978;
Table~\ref{tab:conditional_descstats} in
Appendix~\ref{app:competition}), confirming that the regressions
exploit within-item variation where FL and non-FL tenders are
comparable.

Public procurement accounts for 12--15\% of GDP in OECD countries and
a larger share in developing economies \citep{oecd2019government}, where
bid rigging inflates costs by an estimated 10--50\%
\citep{connor2007price}. Existing detection relies on \emph{reactive}
methods (leniency, whistleblowers) or \emph{proactive} screens that
require granular bid-level data \citep{harrington2008detecting,
imhof2019detecting}. In many developing countries, such data are
unavailable or unreliable. The FL screen addresses this gap: it requires
only participation and outcome records, produces firm-level flags that
implicate all tenders in which a flagged firm participates, and yields
an operationally implementable detection threshold.

This paper advances two hypotheses of increasing strength:

\begin{description}
  \item[H1 (Screening marker):] FL presence is a useful
    \emph{screening marker} for potential collusion---it correlates
    with price anomalies and overlaps with known cartel participants,
    regardless of whether every FL firm is itself a cover bidder.
  \item[H2 (Cover-bidding mechanism):] FL firms are \emph{themselves}
    cover bidders deployed by cartels to simulate competition---FL
    presence marks cartel operations that sustain supra-competitive
    prices.
\end{description}

\noindent Our primary contribution is to establish H1: FL presence as a
robust screening marker for bid rigging. We do not claim to identify
the causal effect of FL presence on prices. The OLS, cross-fit, and
matching estimates capture a \emph{conditional association}---the
average price difference between FL-present and FL-absent tenders
within the same item type, year, and purchasing unit. This association
is consistent with cartels deploying cover bidders in tenders where
they also inflate prices, but it could also reflect unobserved
market-level factors that attract both FL firms and higher prices.
Sensitivity analysis (\citealt{cinelli2020making}, $RV = 17.5\%$)
bounds the scope for omitted-variable bias, and we discuss remaining
threats throughout.

Independent empirical diagnostics---co-bidding network analysis,
bid-level coordination tests, and structural estimation---then provide
suggestive evidence for the cover-bidding mechanism (H2). None of these
tests is individually conclusive; their value lies in the consistency
of the pattern across distinct approaches.

\subsection{Contributions}

We make three contributions, each leveraging the scale and
institutional detail of the BEC platform.

\emph{First}, we propose FL status as a novel, operationally
implementable screening marker for bid rigging. The screen requires
only participation and outcome records---no bid-level data---and
produces firm-level flags for targeted investigation. External
validation against CADE (Brazil's competition authority) convictions
shows that 7.1\% of FL firms co-participate in tenders with convicted
cartelists---3.5 times the rate expected by chance ($p < 0.001$,
permutation test). Co-participation captures market proximity to known
cartels, not direct cartel membership; it indicates that FL firms
systematically operate in cartel-affected markets. The FL--price
association persists after excluding CADE-involved markets
($\hat{\beta} = 0.062$, $N = 1{,}622{,}954$), indicating that the
screen detects anomalies beyond known cartels. The FL screen
complements bid-level tools such as \citet{imhof2019detecting} in a
two-stage workflow: FL screening triages firms using participation
data; bid-level analysis then examines flagged tenders where
granular data are available.

\emph{Second}, we document a robust conditional association between FL
presence and procurement prices across multiple estimation approaches.
OLS with high-dimensional fixed effects yields a baseline of 6.4\%;
temporal cross-fitting yields 3.6\%; CEM and IPW matching yield
5.5--7.7\%. These estimators share the maintained assumption of
selection on observables (plus fixed effects); they differ in
functional form and reweighting method. A leave-one-out IV
diagnostic (19.4\%) suggests the OLS is attenuated by
measurement error in the binary FL indicator, but exclusion-restriction
concerns preclude treating the IV as a primary estimate
(Section~\ref{sec:results_ols}). The consistent 3.6--7.7\% range
across approaches with different functional-form assumptions provides
confidence in the association, while the sensitivity analysis
($RV = 17.5\%$) bounds the residual scope for omitted-variable
bias.

\emph{Third}, we provide suggestive evidence for a cover-bidding
mechanism (H2) from multiple independent diagnostics. Co-bidding
network analysis shows that the FL--price effect concentrates in
competitive markets (low winner HHI), where the structural model
predicts cover bidding is most profitable. \cites{bajari2003deciding}
tests reject exchangeability ($D = 0.15$, $p < 0.001$) and
conditional independence (pairwise product = 4.28, $t = 81.0$) of
FL bid residuals. Structural estimation favors coordinated cover
bidding (Regime~2; $\hat{\sigma}_c / \hat{\sigma}_g = 0.72$). A
framework for cover-bidder deployment generates testable predictions
that organize these diagnostics and informs illustrative policy
counterfactuals.

The remainder of this paper is organized as follows.
Section~\ref{sec:data} describes the data and FL definition.
Section~\ref{sec:structural} develops a framework for cover-bidder
deployment with testable predictions. Section~\ref{sec:empirical}
details the empirical strategy. Section~\ref{sec:cade} validates the
FL screen against CADE cartel convictions. Sections~\ref{sec:results}
and~\ref{sec:mechanisms} present results and mechanism tests.
Section~\ref{sec:robustness} reports robustness checks and
illustrative welfare estimates.
Section~\ref{sec:conclusion} concludes.
